\chapter{Configuration Manifolds}
\label{ch:configuration_manifolds}

\section{Describe the System Before Choosing Coordinates}

A club has a position and an orientation. An elbow angle describes a relative
configuration, while clubface orientation describes an output of the entire
kinematic chain. Angular velocity adds information that a configuration does
not contain. These distinctions matter when comparing swings: a small change
in a joint coordinate can accompany a large change in an impact outcome, and
a large coordinate jump can occur while the physical motion remains smooth.
Geometry helps us distinguish these possibilities. It does not determine which
motion a golfer should make without a mechanical model and an outcome criterion.

Coordinates are measurements expressed as numbers. They are legitimate local
descriptions, not false versions of the physical system. A useful model states
what configurations are identified, what constraints apply, and where its
coordinates remain valid. Some mechanical configuration spaces are Euclidean;
others require multiple charts or a constrained representation. The relevant
question is whether the chosen description covers the motion being analyzed
without ambiguity that the analysis fails to handle.

An ideal revolute joint with unlimited rotation has configuration space $S^1$
when a full turn restores every modeled property. If cable winding, accumulated
twist, or another internal variable distinguishes successive turns, that angle
alone is not a complete configuration. A joint restricted between two stops
instead has an interval of admissible angles. Its interior admits a single
real coordinate; the endpoints require a boundary or contact model. These
models answer different questions and should not be interchanged silently.

A smooth $n$-dimensional manifold is a Hausdorff, second-countable space locally
described by open subsets of $\mathbb R^n$, with smooth invertible transitions
between overlapping coordinate charts. This definition covers $\mathbb R^n$
itself. A chart need not cover the entire space, but some manifolds do admit one
global chart. A circle does not admit a single nonsingular real coordinate
with all its identifications preserved. Embedding it as $(\cos\theta,\sin\theta)$
uses two numbers subject to one constraint rather than creating a second
degree of freedom. Boundaries, corners, changing contact modes, and singular
constraint sets require qualifications beyond an ordinary smooth manifold.

\subsection{Configurations, Directions, and Task Outputs}

A spherical pendulum whose only configuration is the direction of a rigid rod
has space $S^2$. An unconstrained rigid-body orientation has space $SO(3)$,
which also records rotation about the rod's axis. A two-axis universal joint
consists of ordered revolute joints. With ideal unlimited independent angles,
its joint configuration space is $S^1\times S^1$, not automatically $S^2$.
Its pointing-direction output may cover a sphere with repeated or singular
representations. Joint limits, interference, or closure can restrict that
space. Counting two degrees of freedom does not identify its topology.

An open ideal chain with two spherical joints and three revolute joints has
the nine-dimensional configuration space
\begin{equation}\label{eq:golf_config}
 Q_0=SO(3)\times SO(3)\times(S^1)^3,\qquad \dim Q_0=9.
\end{equation}
This can be a teaching model of torso, shoulder, elbow, and wrist motion.
It is not an anatomical assertion that those joints rotate without limits,
nor a complete two-handed golfer. A real model must state its pelvis/ground
attachment, joint axes, ranges, hand contacts, and whether shaft flexibility
or grip compliance is represented. Shoulder translation and distributed
spinal motion may matter for the chosen question.

Smooth holonomic constraints $\Phi(q)=0$ with independent Jacobian rows define
a local submanifold of dimension $n-\operatorname{rank}D\Phi$. Velocities then
satisfy $D\Phi(q)\dot q=0$. At changing rank, this dimension argument needs
reassessment. Unilateral contact also has inequalities and mode transitions.
A velocity constraint need not integrate to a configuration constraint.
Consequently, a closed chain cannot be analyzed by adding nominal joint
degrees of freedom and ignoring closure. Nor does configuration dimension
give the number of independently identifiable hand forces or muscle inputs.

The clubhead position $p(q)$ and face orientation $R_c(q)$ are task maps out
of this configuration space. Their derivatives can lose rank even when the
configuration coordinates are regular. Conversely, singular coordinates can
be replaced without changing the rank of a physical task map. Keeping these
two questions separate is essential when interpreting a large wrist-angle
change or a poorly conditioned inverse calculation.

\section{Velocities, Forces, and Mechanical State}

A tangent vector at $q$ describes an infinitesimal configuration change. In
coordinates it has components $v^i$, and a mechanical velocity is one such
vector. The tangent bundle $TQ$ collects all pairs $(q,v)$. For a regular
second-order mechanical model with no additional internal state, it is a
natural state space of dimension $2n$. Muscle activation, viscoelastic memory,
flexible modes, controller state, and contact mode can require an augmented
state. A configuration path alone does not determine any of these quantities.

A generalized force is a covector $F\in T_q^*Q$: it acts on a virtual
displacement to give virtual work, and on velocity to give power. In coordinates,
\begin{equation}\label{eq:manifold_force_power}
 \delta W=F_i\,\delta q^i,\qquad P=F_i\dot q^i.
\end{equation}
Repeated indices are summed here. A torque represented by a three-vector in
a chosen orthonormal frame uses an identification between vectors and
covectors; this convenience should not obscure the dual role of force.
Mass and inertia provide a different identification in generalized coordinates.

Let new coordinates be $z=\psi(q)$ with invertible Jacobian $A=D\psi$.
Velocity components transform as $v_z=A v_q$, whereas force components
transform as $F_z=A^{-T}F_q$. Therefore $F_z^Tv_z=F_q^Tv_q$: physical power
does not depend on the coordinate units or chart. Transforming both with $A$
would generally destroy this equality. This is a direct diagnostic for
incorrect unit scaling in a swing model.

For a positive-definite mass matrix $M(q)$, kinetic energy is
\begin{equation}\label{eq:manifold_kinetic_metric}
 T(q,v)=\tfrac12v^TM(q)v,\qquad
 M_z=A^{-T}M_qA^{-1}.
\end{equation}
It defines an inner product on configuration tangent vectors. Redundant
coordinates, massless idealizations, or uneliminated constraints can instead
give a singular matrix; then a positive-definite metric must be justified on
the appropriate reduced tangent space. The configuration mass matrix is not
automatically a metric on the augmented position/velocity/activation state.

\section{Rotations and Three Different Singularities}

We use active rotation matrices $R\in SO(3)$ mapping body components into
world components. Thus $R^TR=I$ and $\det R=1$. Body angular velocity $\omega$
and world angular velocity $\omega_w$ satisfy
\begin{equation}\label{eq:manifold_rotation_kinematics}
 \dot R=R\widehat\omega=\widehat{\omega_w}R,\qquad
 \omega_w=R\omega,\qquad \widehat a\,b=a\times b.
\end{equation}
Nine matrix entries describe three rotational degrees of freedom subject to
constraints. A numerical integrator must preserve or appropriately restore
these constraints. This representation avoids Euler-chart singularities,
but it does not remove joint stops, loss of task authority, or input bounds.

For the specific ZYX convention $R=R_z(\psi)R_y(\theta)R_x(\phi)$, body
angular velocity is related to roll, pitch, and yaw rates by
\begin{equation}\label{eq:manifold_euler_rates}
 \omega=
 \begin{pmatrix}
 1&0&-\sin\theta\\
 0&\cos\phi&\sin\phi\cos\theta\\
 0&-\sin\phi&\cos\phi\cos\theta
 \end{pmatrix}
 \begin{pmatrix}\dot\phi\\\dot\theta\\\dot\psi\end{pmatrix},
 \qquad \det E=\cos\theta.
\end{equation}
At $\theta=\pm\pi/2$ the rate map loses rank. Recovering arbitrary angle
rates from a physical angular velocity is then impossible or nonunique, and
nearby inversions can amplify measurement errors. This does not mean the
rigid body loses a rotational degree of freedom. A pure pitch motion
$R(t)=R_y(t)$ crosses $\pi/2$ with finite body velocity $(0,1,0)$; this special
coordinate path remains finite even though the chart cannot represent all
nearby motions regularly. Infinite torque does not follow from the chart
alone: it depends on whether a controller performs an ill-conditioned
inversion, differentiates discontinuities, saturates, or switches charts.

There are three distinct concerns. A \emph{coordinate singularity} is a
failure of a representation. A \emph{mechanical singularity} may align actual
gimbal axes or change the mobility of a mechanism. A \emph{task singularity}
is a rank loss in the map from feasible joint velocity to the selected output.
A three-gimbal inertial platform can physically suffer gimbal alignment;
rewriting its orientation with quaternions does not rebuild its hardware.
Historical spacecraft incidents should not substitute for this distinction.

For biomechanics, report the rotation sequence, segment frames, calibration,
and the observed distance to its singular set. A different Euler sequence or
local chart can be entirely adequate over a measured swing. Quaternion or
matrix methods are useful when coverage, interpolation, or integration makes
them preferable. No sequence-independent claim that every golf shoulder
passes near gimbal lock is justified without actual orientation data.

\subsection{Exponential Coordinates and the Half-Turn Cut}

For a rotation vector $a\in\mathbb R^3$, $\alpha=\|a\|$, Rodrigues' formula is
\begin{equation}\label{eq:manifold_rodrigues}
 \exp(\widehat a)=I+\frac{\sin\alpha}{\alpha}\widehat a
       +\frac{1-\cos\alpha}{\alpha^2}\widehat a^2.
\end{equation}
The coefficients have limits $1$ and $1/2$ at zero. Stable implementations
use appropriate small-angle expansions. The principal rotation logarithm
chooses an angle below $\pi$ away from half-turns. At angle $\pi$, axes $n$
and $-n$ give the same rotation with opposite rotation vectors. A single
continuous principal branch cannot resolve that choice globally.

This branch ambiguity differs from a singular differential of the exponential
map. The right-trivialized differential is
\begin{equation}\label{eq:manifold_exp_jacobian}
 J_r(a)=I-\frac{1-\cos\alpha}{\alpha^2}\widehat a
       +\frac{\alpha-\sin\alpha}{\alpha^3}\widehat a^2,
 \qquad \det J_r=\left(\frac{2\sin(\alpha/2)}{\alpha}\right)^2.
\end{equation}
Its determinant is $4/\pi^2$ at a half-turn and zero at a nonzero full turn.
The exponential is locally regular at each half-turn vector even though the
principal logarithm must choose between globally identified vectors. At
zero the determinant limit is one. This example separates local invertibility
from global uniqueness, two ideas that should remain distinct throughout
trajectory analysis.

A unit quaternion is another rotation representation, with $q$ and $-q$
representing the same orientation. The raw Euclidean difference can have norm
two for identical physical orientations. A shortest angular distance is
$2\arccos(|q_1^Tq_2|)$ for normalized quaternions, with the dot product clipped
for numerical roundoff. Sign-consistent interpolation or a sign-invariant
error is needed. A raw quaternion difference does not generally reverse
direction when only one input sign changes; the actual defect is dependence
on an arbitrary representative. Lifting rotations to quaternions does not
by itself guarantee global continuous feedback convergence or prevent unwinding.

\section{What a Kinetic Metric Measures}

The kinetic metric gives the length of a configuration curve $q(\lambda)$ as
\begin{equation}\label{eq:manifold_metric_length}
 L_M=\int\sqrt{q'(\lambda)^TM(q(\lambda))q'(\lambda)}\,d\lambda
     =\int\sqrt{2T(q(t),\dot q(t))}\,dt.
\end{equation}
It is invariant under regular monotone changes of the path parameter. It is
not an energy expenditure. With ordinary mechanical units its length has
units $\sqrt{\mathrm{kg}}\,\mathrm m$, whereas its time derivative has units
$\sqrt{\mathrm J}$. Work also depends on applied forces, and metabolic cost
requires a physiological model. The action $\int T\,dt$ is another quantity
again and depends on timing.

For an illustrative pair of independent constant-inertia rotations with
$I_1=12$ and $I_2=0.03\;\mathrm{kg\,m^2}$, equal angular displacements have
metric lengths in ratio $\sqrt{I_1/I_2}=20$. At equal angular speeds their
kinetic energies have ratio $400$. If the first motion takes twenty times
as long with a correspondingly scaled speed profile, their instantaneous
kinetic energies at matched path phase can instead be equal. The angle
change alone establishes neither energy ratio nor required work. These
numbers illustrate a model; they are not subject-specific torso and wrist
inertia estimates. Coupled limbs also have off-diagonal mass-matrix terms.

Different metrics answer different questions. An orientation distance based
on rotation angle is convenient for face alignment; a measurement covariance
metric can quantify uncertainty; the kinetic metric weights generalized
velocity by inertia. For the standard rotation-angle metric, the largest
distance on $SO(3)$ is $\pi$. Arbitrary scaling or an anisotropic kinetic
metric changes that statement. A clubface error measured in degrees should
not be silently reinterpreted as a kinetic distance or a success probability.

\section{Geodesics, Gravity, and Passive Dynamics}

The Levi-Civita connection of a metric is the unique metric-compatible,
torsion-free connection. Its coordinate coefficients are
\begin{equation}\label{eq:manifold_christoffel}
 \Gamma^i_{jk}=\tfrac12(M^{-1})^{i\ell}
 (\partial_jM_{\ell k}+\partial_kM_{\ell j}-\partial_\ell M_{jk}),
 \qquad
 (\nabla_{\dot q}\dot q)^i=\ddot q^i+\Gamma^i_{jk}\dot q^j\dot q^k.
\end{equation}
An affinely parameterized geodesic satisfies $\nabla_{\dot q}\dot q=0$.
It has constant metric speed and minimizes length on sufficiently short
segments. A long geodesic need not be a globally shortest path. General
reparameterization preserves its unparameterized image but can introduce
tangential covariant acceleration. The geodesic equation requires its
parameter convention as well as its metric.

Let $F=B(q)u+F_{\rm passive}+F_{\rm external}$ denote nonpotential generalized
force covectors, and let $V$ include the conservative gravity and elastic
potentials chosen for the model. Avoid counting an elastic force in both
$V$ and $F_{\rm passive}$. For an unconstrained regular simple mechanical
system, the equivalent covector and vector equations are
\begin{equation}\label{eq:euler_lagrange_manifold}
 M\nabla_{\dot q}\dot q=-dV+F,\qquad
 \nabla_{\dot q}\dot q=-\operatorname{grad}_M V+M^{-1}F,
 \qquad \operatorname{grad}_M V=M^{-1}dV.
\end{equation}
Here $dV$ means the column of partial derivatives when written in coordinates.
The metric lowers a tangent index; its inverse raises a covector index.
Multiplying the left side by $M$ while leaving a Riemannian gradient on the
right would mix the two types. In coordinates the same mechanics is
$M\ddot q+C\dot q+dV=F$, where $C\dot q$ equals the metric-lowered connection
term. The matrix $C$ is not uniquely determined by that product.

Only when the total nonkinetic force vanishes does physical time parameterize
a geodesic of the kinetic metric. Zero actuator command does not eliminate
gravity, springs, dissipation, or external forces. An ideal pendulum with
$M=m\ell^2$ has zero connection coefficient in its angular coordinate, yet
under gravity it satisfies $\ddot\theta=-(g/\ell)\sin\theta$. With $\ell=1$
metre and $\theta=\pi/6$, its acceleration is $-4.905\;\mathrm{rad/s^2}$.
It is passive and conserves $T+V$, but is not a time-parameterized geodesic
of that kinetic metric. Fixed-energy conservative paths can sometimes be
related to a Jacobi metric with a different parameter; this is a separate
construction, requiring $E>V$ and excluding the degeneracy at turning points.

Thus an observed follow-through cannot be called nearly geodesic merely
because it looks relaxed. The test is a model-based covariant-acceleration
residual, interpreted alongside estimated gravity, elasticity, contact,
dissipation, and active moments. A small residual is not an observation of
small muscle activation. Opposing muscles can have substantial activation
and small net moment, and motion data may not identify those forces uniquely.

\subsection{A Flat-Space Counterexample}

In polar coordinates on the Euclidean plane away from the origin, the unit-mass
metric is $M=\operatorname{diag}(1,r^2)$. Its nonzero coefficients include
$\Gamma^r_{\theta\theta}=-r$ and
$\Gamma^\theta_{r\theta}=\Gamma^\theta_{\theta r}=1/r$. A free particle obeys
\begin{equation}\label{eq:manifold_polar_geodesic}
 \ddot r-r\dot\theta^2=0,\qquad
 \ddot\theta+2\dot r\dot\theta/r=0.
\end{equation}
For the straight Cartesian motion $(x,y)=(1,t)$, take
$r=\sqrt{1+t^2}$ and $\theta=\arctan t$. Then
$\dot r=t/r$, $\ddot r=1/r^3$, $\dot\theta=1/r^2$, and
$\ddot\theta=-2t/r^4$. Both equations vanish exactly, despite nonzero
ordinary coordinate accelerations and connection coefficients.

Intrinsic Riemann curvature is zero here. Connection coefficients are not
themselves curvature tensors; they can be nonzero in curvilinear coordinates
on flat space. Conversely, the topology of a configuration space does not
specify its metric curvature. The coordinate terms represent a consistent
description of inertial dynamics, not a source of energy. In a multibody
system, inertial coupling can exchange energy between motions while the
overall power balance remains
\begin{equation}\label{eq:manifold_energy_balance}
 \frac{d}{dt}(T+V)=F_i\dot q^i
\end{equation}
for time-independent $M,V$ and the stated unconstrained model. Stationary
ideal constraint reactions do no work on allowed velocities; moving constraints,
friction, and impacts require their own power or energy terms. A distal speed
increase must be reconciled with this balance, not attributed to work created
by a coordinate connection.

\section{Path Bending and Required Actuation}

Let $q(s)$ be a regular configuration path parameterized by kinetic-metric
arc length, $T_q=dq/ds$ its unit tangent, and $\nu=ds/dt>0$ its metric speed.
To avoid confusing kinetic energy with the tangent, the subscript is retained.
Its covariant curvature vector is $k_g=\nabla_{T_q}T_q$, with norm
$\kappa_g=\|k_g\|_M$. The acceleration decomposition is
\begin{equation}\label{eq:manifold_path_force}
 \nabla_{\dot q}\dot q=\dot\nu T_q+\nu^2 k_g,\qquad
 B u=M(\dot\nu T_q+\nu^2 k_g)+dV-F_{\rm passive}-F_{\rm external}.
\end{equation}
The curvature vector is metric-orthogonal to the tangent. In the special
case that actuators supply the entire force and there is no potential or
other load, the transverse force covector is $\nu^2 M k_g$, whose dual norm
is $\nu^2\kappa_g$. The dual norm is $\|F\|_{M^{-1}}=\sqrt{F^TM^{-1}F}$.
This is not an unqualified formula for muscle effort or Cartesian newtons.

In a constrained system add reaction covectors and enforce the constraints,
or use a valid reduced model. In an underactuated system the right side must
lie in the image of $B$, and a feasible solution must satisfy input bounds.
Different columns can have different torque limits and biological costs.
The scalar curvature does not establish this feasibility. A guide can provide
normal reaction without actuator work; gravity can supply some acceleration
while changing potential energy. The feasible timing calculation in Chapter 2
is therefore still required.

Geodesic curvature is not obtained by subtracting a scalar coordinate curvature
from a scalar ``Coriolis curvature.'' It is a norm of a covariant derivative
in a specified metric. Ordinary curvature of a plotted coordinate path changes
under coordinate scaling and can have different units. Neither is automatically
the Cartesian curvature of the clubhead path $p(q)$. That physical path still
obeys $\ddot p=J\ddot q+\dot J\dot q$. These are connected descriptions,
but they answer different questions about the same motion.

\section{Tracking Errors With Consistent Frames}

A nearby configuration can be compared with a reference using a common chart,
a local Riemannian logarithm, a retraction, or an embedded error with stated
constraints. A Lie-group logarithm is another option when the configuration
space actually has that group structure. It is not a universal subtraction
operation on every manifold. The Riemannian exponential of a chosen metric
also need not coincide with the Lie-group exponential; they agree in special
settings such as the standard bi-invariant rotation metric.

Velocities at different configurations require a stated identification.
Levi-Civita parallel transport is one choice and depends on the connection
and path. A common coordinate chart or group trivialization is another.
Rotating components between body frames is a physical frame transformation;
it should not be casually labeled parallel transport for an arbitrary kinetic
metric. These conventions must be settled before constructing feedback gains.

\subsection{A Derived Rigid-Body Tracking Law}

Consider a fully actuated rigid body with constant body inertia $J>0$,
$J\dot\omega+\omega\times J\omega=\tau$. Assume an exactly known model,
a smooth feasible desired rotation $R_d(t)$, and available torque without
saturation for the following derivation. Define $\dot R_d=R_d\widehat\omega_d$
and $A=R^TR_d$, which maps desired-body components into actual-body components.
Let
\begin{equation}\label{eq:manifold_attitude_errors}
 e_\omega=\omega-A\omega_d,\qquad
 \Psi=\tfrac12\operatorname{tr}(I-R_d^TR),\qquad
 e_R=\tfrac12(R_d^TR-R^TR_d)^\vee.
\end{equation}
The vee map inverts the hat map on skew matrices. This smooth $e_R$ is
$\sin\alpha$ times the relative rotation axis, not the logarithmic angle-axis
vector $\alpha n$. In particular it vanishes at a half-turn, an important
limitation rather than a numerical defect to conceal.

Differentiating the frame transformation gives
\begin{equation}\label{eq:manifold_attitude_transport}
 \dot A=-\widehat\omega A+A\widehat\omega_d,\qquad
 b=\frac{d}{dt}(A\omega_d)=-\omega\times(A\omega_d)+A\dot\omega_d.
\end{equation}
Using the inverse transformation $R_d^TR$ in $e_\omega$ would compare vectors
in inconsistent frames. Omitting the first term in $b$ would omit the rate
at which the reference components rotate relative to the actual body.

For positive scalar gains, choose
\begin{equation}\label{eq:pd_so3}
 \tau=-k_R e_R-k_\omega e_\omega+\omega\times J\omega+Jb.
\end{equation}
Substitution yields $J\dot e_\omega=-k_R e_R-k_\omega e_\omega$.
Direct differentiation of $\Psi$ gives $\dot\Psi=e_R^Te_\omega$; hence
\begin{equation}\label{eq:manifold_attitude_energy}
 W=\tfrac12e_\omega^TJe_\omega+k_R\Psi,\qquad
 \dot W=-k_\omega\|e_\omega\|^2.
\end{equation}
For $W(0)<2k_R$, the trajectory stays away from the half-turn potential level.
In the ideal closed error dynamics, the largest invariant zero-dissipation
set in this region has $e_\omega=0$, then $e_R=0$, and therefore the desired
relative orientation. This supplies a local attraction argument. Stronger
rate or domain claims require the corresponding proof and assumptions.
At a half-turn with zero velocity error the feedback vanishes, providing an
explicit counterexample to global attraction. A globally defined smooth
formula is not the same as a globally convergent controller. Logarithmic
feedback has its own branch restriction; quaternion feedback must respect
the double cover, and hybrid designs require analysis of their switching.

Two checks make the formula practical to audit. If actual and desired bodies
have the same world angular velocity $w$, then $\omega=R^Tw$ and
$\omega_d=R_d^Tw$, so $e_\omega=0$ exactly. For arbitrary orientations and
rates, substituting the torque into the rigid-body equation must reproduce
the stated $\dot W$. Both identities can be verified numerically without
simulating a golf swing. Torque saturation, disturbances, uncertain inertia,
underactuation, and muscle activation dynamics change this closed-loop result.
It is a mathematical example, not an identified human motor-control law.

\section{Connecting Geometry to a Golf Investigation}

Begin with the observable outcome: preimpact clubhead velocity, face orientation,
strike location, and their uncertainty at a defined event. Construct the
kinematic model needed to map measured segment poses to those outcomes.
State which joints, constraints, flexible elements, and contact modes the
model retains. A larger geometric model is useful only when its parameters
and outputs can be estimated well enough for the question.

Next verify representation consistency. Record active/passive rotation
conventions, handedness, body/world frames, joint sequence, units, quaternion
ordering, and calibration. Test that converting coordinates preserves physical
orientations, velocities, kinetic energy, and power. Examine chart conditioning
along the actual data. A discontinuity created by angle wrapping is not
evidence of a sudden physical release, and smoothing it as if it were physical
can corrupt estimated acceleration and inverse torque.

Then connect kinematics to dynamics. Estimate the forces and moments required
by the motion, including gravity, elastic energy, hand and ground interactions,
and the shaft model. Propagate measurement and parameter uncertainty. Compute
covariant quantities only with a declared metric and compare them with physical
clubhead quantities through the task Jacobian. Apparent simplicity in one
coordinate system does not establish lower loading, greater efficiency, or
better impact reliability.

Finally test the explanatory hypothesis. If inertial coupling is proposed to
assist distal speed, verify the energy and power balance and compare forward
simulations under a precisely defined change in coupling, activation, or timing.
Recompute the resulting motion rather than subtracting a term from a fixed
trajectory and calling it a performance prediction. Check predictions on
held-out swings and report where uncertainty prevents discrimination between
mechanisms. Passive assistance, active regulation, and task constraints may
operate together; their contributions are questions to estimate.

A configuration tube $d_M(q,q_d(t))\leq\varepsilon$ bounds only configuration
error under its chosen metric. Two states inside it can have very different
velocities and impact results. A state tube needs a metric or certificate on
the actual state space, with explicit velocity comparisons and contact events.
Invariance must be verified under the closed-loop dynamics and disturbances.
Geometry defines the set; dynamics determine whether a trajectory stays inside,
and the task map determines whether staying inside is useful.

The benefit of manifold methods is precise bookkeeping of configuration,
velocity, force, and frame relationships across a motion. That precision
strengthens an explanation of the swing when combined with measured evidence.
It does not remove the need for timing, feasibility, stability, or an empirical
account of how the golfer supplies and regulates forces.

\section{Further Reading}

\href{https://hades.mech.northwestern.edu/images/2/25/MR-v2.pdf}{Lynch and Park's
Modern Robotics} provides joint, rotation, wrench, and rigid-body conventions.
\href{https://arxiv.org/abs/1003.2005}{Lee, Leok, and McClamroch's geometric
control analysis} develops smooth attitude tracking with explicit attraction
conditions. These are mathematical and engineering sources; applying their
structures to a human swing requires the additional modeling and validation
steps described here.

\section{Exercises}

\begin{enumerate}[itemsep=3pt,parsep=0pt]
\item Compare an unlimited revolute joint, a joint with stops, and a rod whose
direction alone is measured. Identify the configuration, output map, and any
unobserved rotation. Explain why dimension alone cannot identify topology.
\item For an ideal chain with two spherical and three revolute joints, count
configuration and tangent-bundle dimensions. Add two independent holonomic
constraints locally. Explain why an activation state changes the state count.
\item Derive the ZYX body-rate matrix and its determinant. Evaluate a pure
pitch crossing through the rate-map singular set at $\theta=\pi/2$.
Distinguish loss of inverse uniqueness from unbounded physical velocity.
\item Evaluate $\det J_r$ at $0$, $\pi$, and $2\pi$. Compare the exponentials
of $\pi n$ and $-\pi n$. Explain why local invertibility at each half-turn
does not define a globally continuous principal logarithm.
\item Choose a nonsingular coordinate scaling and verify the velocity, force,
and mass-matrix transformation laws. Check kinetic energy, power, and dual
force norm numerically before and after transformation.
\item Verify the polar-coordinate free-particle example and compute a component
of its Riemann curvature tensor. Contrast nonzero connection coefficients
with zero intrinsic curvature and zero Cartesian acceleration.
\item Derive the gravity pendulum equation and its energy balance. Compare
its motion with the kinetic-metric geodesic from the same initial state.
Explain what zero actuator torque does and does not imply.
\item For a regular configuration path, derive the required actuator covector
including a potential and external force. State the image and bound tests
needed when the actuation map is rectangular.
\item Implement the attitude law with two different body orientations but a
common world angular velocity. Verify zero velocity error, the transport
derivative, and the dissipation identity. Test the undesired half-turn
equilibrium and a torque-saturated case separately.
\item Design a swing-data comparison that preserves the physical motion while
changing its coordinates. Identify quantities that must remain invariant and
quantities that depend on the chosen metric. Add a separate intervention that
changes the dynamics and specify the impact prediction it would test.
\end{enumerate}
